% PAGES: 250-250

Then, since $Se^{-qT} = PVF$, we obtain that
\begin{equation}
f_P'(x) \;=\; PVF\sqrt{\frac{T}{2\pi}}\,\exp\!\left(-\frac{(d_1(x))^2}{2}\right).
\end{equation}

Note that $f_P'(x) = f_C'(x)$, since $vega_P = vega_C$.

The Newton's method recursion for solving (8.45) is
\begin{equation}
x_{k+1} \;=\; x_k \;-\; \frac{f_P(x_k)}{f_P'(x_k)},
\end{equation}
where the functions $f_P(x)$ and $f_P'(x)$ are given by (8.46) and (8.48),
respectively.

A good initial guess for Newton's method is 25\% volatility, i.e., $x_0 = 0.25$, and the
algorithm is stopped when two consecutive approximations in Newton's method are within
$10^{-6}$ of each other; see the pseudocode from Table 8.3 for finding the implied
volatility for both call and put options, i.e., for solving either (8.44) or (8.49).

\begin{table}[H]
\centering
\caption{Pseudocode for computing implied volatility}
\fbox{\begin{minipage}{0.85\linewidth}
Input:\\
$V_m = $ option price\\
\hspace*{1.5em}\textit{// $V_m = C_m$ for call implied vol; $V_m = P_m$ for put
implied vol}\\
$K = $ strike price of the option\\
$T = $ maturity of the option\\
$PVF = $ present value of the forward price of the underlying asset\\
$disc = $ discount factor corresponding to time $T$\\
$tol = $ tolerance for Newton's method convergence\\
$f_{BS}(x) = $ Black--Scholes option value; $x = $ volatility\\
\hspace*{1.5em}\textit{// $f_{BS}(x) = f_C(x)$ for calls; $f_{BS}(x) = f_P(x)$ for
puts}\\
\hspace*{1.5em}\textit{// $f_{BS}'(x) = f_C'(x) = f_P'(x)$}
\bigskip

Output:\\
$x_{new} = $ implied volatility
\bigskip

$x_0 = 0.25;$ \hfill\textit{// initial guess: 25\% volatility}\\
$x_{new} = x_0; x_{old} = x_0 - 1; tol = 10^{-6}$\\
while ($|x_{new} - x_{old}| > tol$)\\
\hspace*{1.5em}$x_{old} = x_{new}$\\
\hspace*{1.5em}$x_{new} = x_{old} - \dfrac{f_{BS}(x_{old}) - V_m}{f_{BS}'(x_{old})}$\\
end
\end{minipage}}
\end{table}

For the options from Table 8.2, note that the options maturity is $T = \frac{199}{252}$,
i.e., the ratio of 199, the number of trading days between March 9, 2012, and December
22, 2012, and 252, the total number of trading days in a year. Using the values
$PVF = 1349.54$ and $disc = 0.9964$ computed using least squares and the Newton's
method from Table 8.3, we obtain the implied volatilities from Table 8.4.

A consequence of the Put--Call parity is that the theoretical values of implied
volatilities of calls and puts with the same strike are equal. Note that, indeed, the
implied volatilities from Table 8.4 corresponding to calls and puts with the same
strike are nearly identical.
